\section{The marked cubic-surface contraction and the toric hexagon}
\label{sec:picard-hexagon}
This section reconstructs the explicit geometric comparison in the supplied
discussion, lines 2017--2055. All point labels and directions below are part
of the objects. The split degree-six model and its six boundary curves occur
in \href{https://www.math.uni-bielefeld.de/lag/man/290.pdf}{Mark Blunk,
\emph{Del Pezzo surfaces of degree 6 over an arbitrary field}, \S2,
Proposition 2.1, printed p.~3}; the Picard maps, signs and torus coordinates
needed for the present comparison are proved here.

\subsection{The marked surfaces, contraction and complete inverse fibres}
Work over $\mathbb C$. Let $p_1,\ldots,p_6$ be six distinct marked points
of $\mathbb P^2$, no three collinear and not all six on a conic. Choose
coordinates with
\[
 p_1=[1:0:0],\quad p_2=[0:1:0],\quad p_3=[0:0:1].
\]
This is a nonempty explicitly realized domain: one choice is
\[
 p_4=[1:1:1],\qquad p_5=[1:2:3],\qquad p_6=[1:3:2].             \tag{PH1}
\]
For this choice the determinants of the twenty point triples, in
lexicographic index order, are
\[
1,1,3,2,-1,-2,-3,1,-1,-5,1,1,1,-2,-1,1,1,2,1,-3,
\]
with the points as rows in increasing label order; in particular they are all
nonzero. A conic through the first three points has the form
$A x_0x_1+B x_0x_2+C x_1x_2=0$. Its values at the last three have coefficient
matrix
\[
\begin{pmatrix}1&1&1\\2&3&6\\3&2&6\end{pmatrix},\qquad\det=7,
\]
so no conic passes through all six.

Put
\[
Y=\operatorname{Bl}_{p_1,p_2,p_3}\mathbb P^2,
\qquad S=\operatorname{Bl}_{p_1,\ldots,p_6}\mathbb P^2.
\]
Let $q_a\in Y$ be the unchanged point above $p_a$, for $a=4,5,6$.
All three lie in the complement of the coordinate triangle and its
exceptional divisors. The morphism
\[
 f:S=\operatorname{Bl}_{q_4,q_5,q_6}Y\longrightarrow Y             \tag{PH2}
\]
contracts exactly the three exceptional curves $E_4,E_5,E_6$.
Denote the line pullback and exceptional classes on $S$ by $H,E_1,\ldots,E_6$,
and those on $Y$ by $h,e_1,e_2,e_3$.

Here are the fibres and their inverse data explicitly. Near a centre $q_a$
take regular local coordinates $(\xi,\eta)$ vanishing there. The blowup is
\[
 \{(\xi,\eta,[s:t]):\xi t=\eta s\}.
\]
Its two charts are $(\xi,t)$ with $\eta=\xi t$, and $(s,\eta)$ with
$\xi=s\eta$. Over the centre the fibre is the whole $\mathbb P^1$ of
directions $[s:t]$; elsewhere the unique direction is $[\xi:\eta]$.
Thus $f$ is an isomorphism away from the three specified fibres, and
retaining the centres and these exceptional directions gives the full
blowup return. No inverse at a contracted point discards its direction.

The Picard groups and intersection forms are
\[
\operatorname{Pic}(S)=\mathbb ZH\oplus\bigoplus_{i=1}^6\mathbb ZE_i,
\quad H^2=1,\quad H E_i=0,\quad E_iE_j=-\delta_{ij},
\]
and the analogous basis $(h,e_1,e_2,e_3)$ on $Y$. To see the additional
summand at a blowup, push a divisor to the smooth base and subtract its
total transform; the difference is supported on the exceptional curve.
Its coefficient is unique since that curve has self-intersection $-1$,
as follows from the transition of its normal coordinate between the two
charts above. Applying this argument successively gives the displayed bases.
The maps are exactly
\[
 f^*h=H,\quad f^*e_i=E_i\ (1\le i\le3),\qquad
 f_*H=h,\quad f_*E_i=e_i\ (i\le3),\quad f_*E_a=0\ (a\ge4).
                                                               \tag{PH3}
\]
Hence $f_*f^*=1$, the kernel of $f_*$ is
$\mathbb ZE_4\oplus\mathbb ZE_5\oplus\mathbb ZE_6$, and it is orthogonal
to $f^*\operatorname{Pic}(Y)$. For
$D=dH-\sum m_iE_i$, $D'=d'H-\sum m'_iE_i$, one has
\[
 f^*f_*D=D+\sum_{a=4}^6m_aE_a,\qquad
 (f_*D)(f_*D')=DD'+\sum_{a=4}^6m_am'_a.                         \tag{PH4}
\]
These identities prove both the precise lost sublattice and the correction
to intersection numbers. The canonical classes obey
\[
 -K_S=3H-\sum_{i=1}^6E_i=f^*(-K_Y)-E_4-E_5-E_6,
 \quad -K_Y=3h-e_1-e_2-e_3,
 \quad K_S^2=3,\quad K_Y^2=6.                                  \tag{PH5}
\]
The canonical-class formula follows in a blowup chart from
$d\xi\wedge d\eta=\xi\,d\xi\wedge dt$ and then from the canonical
class $-3h$ of the plane.

For completeness, these marked blowups give the stated cubic model.
The cubic forms through the six points have dimension four: their six
evaluation conditions are independent, since a conic through any five
points, multiplied by a line avoiding the sixth, separates that sixth
evaluation. Such a conic exists, is unique by Bezout, and is irreducible
because a pair of lines contains at most four of the points. The cubics
have no further base point: for a point outside the six, the pairs whose
joining line contains it form a matching of at most three edges. There are
fifteen perfect matchings of six labels, and each forbidden edge belongs
to three; hence at least six perfect matchings avoid all forbidden edges.
Such a matching supplies a product of three joining lines
nonzero there. At each marked point, matchings with two different partners
give independent first-order directions. Thus the induced anticanonical
system is base-point free on $S$.
It is ample as well. For a nonexceptional irreducible curve of degree $d$
and multiplicities $m_i$, distinct from the six conics above, Bezout with
each such conic gives $2d\ge\sum_{j\ne i}m_j$. Summing gives
$\sum m_i\le12d/5<3d$, so $(-K_S)C>0$. Each of the six conics and every
$E_i$ has anticanonical degree one. Together with $K_S^2=3$, the
Nakai--Moishezon criterion proves ampleness.
The resulting morphism to $\mathbb P^3$ is finite, and its nondegenerate
surface image has degree at least two. The equality
$3=\deg(\text{morphism})\deg(\text{image})$ makes it birational onto a
cubic surface. A general hyperplane pullback is a smooth curve of genus
one by Bertini and adjunction. Its plane cubic image has arithmetic genus
one, so it is smooth; hence the cubic surface has no singular curve.
A hypersurface with no codimension-one singularity is normal. The finite
birational morphism to this normal surface is an isomorphism. In particular
the anticanonical model is a smooth cubic surface.

\subsection{All twenty-seven line classes and their contracted images}
The complete list on $S$ is
\[
 E_i\ (1\le i\le6),\qquad L_{ij}=H-E_i-E_j\ (i<j),\qquad
 Q_i=2H-\sum_{j\ne i}E_j\ (1\le i\le6).                        \tag{PH6}
\]
The last two classes are the strict transforms of the actual joining lines
and the unique conics through the indicated five points. Their curves are
smooth rational, have square $-1$, and anticanonical degree one, so their
images in the cubic model are lines. Conversely a line is a smooth rational
curve with square $-1$ by adjunction. Unless exceptional, it has degree
$d\ge1$ and nonnegative multiplicities satisfying
\[
 \sum m_i=3d-1,\qquad \sum m_i^2=d^2+1.
\]
Cauchy gives $(3d-1)^2\le6(d^2+1)$, hence $d\le2$. At $d=1$ exactly two
multiplicities are one; at $d=2$ exactly five are one. This proves (PH6)
is exhaustive without omitting a line class.

Let $J=\{1,2,3\}$ and $K=\{4,5,6\}$. The entire contraction correspondence
is the following table. A zero divisor class in its second row records a
contracted curve, not an inverse point of a divisor map.
\[
\begin{array}{c|c|c}
\text{curve on }S&\text{image class on }Y&\text{image curve square}\\\hline
E_i\ (i\in J)&e_i&-1\\
E_a\ (a\in K)&0\quad(\text{point }q_a)&\text{not a curve}\\
L_{ij}\ (i,j\in J)&h-e_i-e_j&-1\\
L_{ia}\ (i\in J,a\in K)&h-e_i&0\\
L_{ab}\ (a,b\in K)&h&1\\
Q_i\ (i\in J)&2h-\sum_{j\in J\setminus\{i\}}e_j&2\\
Q_a\ (a\in K)&2h-e_1-e_2-e_3&1
\end{array}                                                     \tag{PH7}
\]
Each noncontracted entry is the actual line or conic with its remaining
marked incidence conditions; coincident divisor classes do not identify
distinct image curves. The total-transform inverse is
\[
 f^*f_*L_{ij}=L_{ij}+\sum_{a\in K\cap\{i,j\}}E_a,
 \qquad f^*f_*Q_i=Q_i+\sum_{a\in K\setminus\{i\}}E_a,           \tag{PH8}
\]
because the original plane curves pass simply through exactly those centres.
Equations (PH3)--(PH4) prove the table and every square in it.

The six curves
\[
 E_1,L_{12},E_2,L_{23},E_3,L_{13}                                \tag{PH9}
\]
avoid the three contraction centres and return isomorphically to their six
namesakes on $Y$. Each has square $-1$. The intersection rules give
$E_iL_{jk}=1$ exactly when $i\in\{j,k\}$; different $E_i$ are disjoint,
and the three displayed $L_{ij}$ are pairwise disjoint. Thus (PH9) is
exactly the cycle with neighbouring intersection one and every other
intersection zero. Its sum is $-K_Y$.

\subsection{A six-cycle realized by Cremona and a marked configuration map}
Put $\alpha=H-E_1-E_2-E_3$, so $\alpha^2=-2$ and $K_S\alpha=0$.
Define the reflection and permutation
\[
 C(D)=D+(D\alpha)\alpha,
 \qquad \sigma=(1\ 3\ 2),\quad
 \sigma(H)=H,\quad\sigma(E_i)=E_{\sigma(i)},
\]
where $\sigma$ fixes the labels in $K$. Their composite is $T=\sigma C$.
The two operations commute, $C^2=1$, and $\sigma^3=1$. Expansion using
$\alpha^2=-2$ proves $C(D)C(D')=DD'$, and both operations fix $K_S$.
Consequently $T$ is an integral intersection isometry, fixes the canonical
class, and $T^3=C$, $T^6=1$. Its exact basis action is
\[
\begin{split}
T(H)&=2H-E_1-E_2-E_3,\\
T(E_1)&=L_{12},\quad T(E_2)=L_{23},\quad T(E_3)=L_{13},\qquad
T(E_a)=E_a\ (a\in K).
\end{split}                                                     \tag{PH10}
\]
In particular it sends the six entries of (PH9) successively to the next,
including $L_{13}\mapsto E_1$, and its order is exactly six.
Its full permutation of all twenty-seven classes is
\[
\begin{gathered}
(E_1\ L_{12}\ E_2\ L_{23}\ E_3\ L_{13}),\qquad
(E_4)(E_5)(E_6),\\
(L_{14}\ L_{34}\ L_{24})\ (L_{15}\ L_{35}\ L_{25})\
(L_{16}\ L_{36}\ L_{26}),\qquad (Q_1\ Q_3\ Q_2),\\
(L_{45}\ Q_6)\ (L_{46}\ Q_5)\ (L_{56}\ Q_4).
\end{gathered}                                                   \tag{PH11}
\]
To verify completeness directly, a mixed class $L_{ia}$ has $L_{ia}\alpha=0$,
so maps to $L_{\sigma(i),a}$; a class $L_{ab}$ in $K$ has product one with
$\alpha$, so maps to $Q_c$ for $K\setminus\{a,b\}=\{c\}$; the latter
has product $-1$ and maps back. A $Q_i$ with $i\in J$ has product zero
and maps to $Q_{\sigma(i)}$. Together with (PH10), these calculations prove
every entry of (PH11); intersection preservation holds for every pair,
not only for neighbouring sides of the hexagon.

This is realized geometrically on the degree-six surface. The quadratic
plane Cremona map and the coordinate permutation are
\[
 \kappa[x_0:x_1:x_2]=[x_1x_2:x_0x_2:x_0x_1],\qquad
 \rho[x_0:x_1:x_2]=[x_1:x_2:x_0].
\]
Their composite is
\[
 F=\rho\kappa:[x_0:x_1:x_2]\dashrightarrow
 [x_0x_2:x_0x_1:x_1x_2].                                       \tag{PH12}
\]
The blowup $Y$ is also the graph closure
\[
 \{([x],[y])\in\mathbb P^2\times\mathbb P^2:
                 x_0y_0=x_1y_1=x_2y_2\}.
\]
Away from the three coordinate points the second coordinate is the
displayed Cremona value; at each point the two blowup charts give its
exceptional direction. Thus switching the two projective factors extends
$\kappa$ to an automorphism of $Y$, and simultaneous permutation of the
coordinates extends $\rho$. Their action on curve classes is the
\emph{pushforward} $T=\sigma C$; the pullback is $T^{-1}=C\sigma^{-1}$.
For instance a general line has Cremona transform a conic through the three
centres, and $E_i$ maps to the opposite joining line, proving (PH10) directly.

On a six-point blowup the exact geometric statement retains the changed
centres. Set $q'_a=F(q_a)$ and
$S'=\operatorname{Bl}_{q'_4,q'_5,q'_6}Y$. Then $F$ lifts to an isomorphism
$\widetilde F:S\to S'$ making $f'\widetilde F=Ff$, with the exceptional
label $a$ sent to the exceptional label $a$ at $q'_a$. Its inverse is the
lift of $F^{-1}$. On exceptional directions its map is the projectivized
invertible derivative $\mathbb P(dF_{q_a})$, obtained by substituting the
two local blowup charts. For (PH1), the new plane centres are exactly
\[
 p'_4=[1:1:1],\qquad p'_5=[3:2:6],\qquad p'_6=[2:3:6].
\]
Thus (PH10)--(PH11) are the complete marked pushforward between these
configurations; they do not assert an automorphism of the fixed generic
six-point configuration. The contraction maps commute with $T$ on their
Picard groups because the kernel spanned by $E_4,E_5,E_6$ is fixed.

\subsection{Torus coordinates, the exact fan, and the cusp-side convention}
Use the following specified coordinates on the dense torus of $Y$:
\[
 u=x_2/x_0,\qquad v=x_2/x_1,
 \qquad [x_0:x_1:x_2]=[v:u:uv].                                \tag{PH13}
\]
Substitution in (PH12) gives the genuine torus automorphism
\[
 F(u,v)=(u/v,u),\qquad F^{-1}(u,v)=(v,v/u),\qquad
 F^3(u,v)=(u^{-1},v^{-1}),\quad F^6(u,v)=(u,v).                  \tag{PH14}
\]
For example $F^2(u,v)=(v^{-1},u/v)$, proving the third-power formula by
one more substitution. At a retained centre $(u,v)$ the exceptional
direction map in (PH2) is completely explicit:
\[
dF_{(u,v)}=\begin{pmatrix}1/v&-u/v^2\\1&0\end{pmatrix},\quad
\det dF=u/v^2,\qquad
[P:Q]\longmapsto[vP-uQ:v^2P].                                  \tag{PH14a}
\]
At the image centre $(U,V)=(u/v,u)$ its inverse sends
$[P':Q']$ to $[U^2Q':UQ'-VP']$. Substitution gives
$[u^2P:u^2Q]=[P:Q]$, proving the inverse on the entire exceptional
$\mathbb P^1$, including vertical and infinite slope directions.
The valuation vectors
$(\operatorname{ord}_B u,\operatorname{ord}_B v)$ of the boundary curves
in the retained order are precisely
\[
\begin{array}{c|rrrrrr}
B&E_1&L_{12}&E_2&L_{23}&E_3&L_{13}\\\hline
\operatorname{ord}_B u&1&1&0&-1&-1&0\\
\operatorname{ord}_B v&0&1&1&0&-1&-1
\end{array}                                                     \tag{PH15}
\]
For a joining line these are its zero and pole orders in (PH13). At $E_1$,
both $x_1$ and $x_2$ vanish to first order while $x_0$ is a unit, giving
$(1,0)$; the other two exceptional rows follow in the same way from the
displayed coordinate ratios. Thus the six rays, in the exact curve order,
are
\[
n_1=(1,0),\ n_2=(1,1),\ n_3=(0,1),\ n_4=(-1,0),\
n_5=(-1,-1),\ n_6=(0,-1).
\]
Adjacent pairs have determinant one. The cones between them are the six
smooth toric charts of the three-point blowup: the original plane fan's
three cones are each subdivided by the sum of their two generating rays,
which is the blowup operation in the local charts of (PH2).
On this cocharacter lattice (PH14) acts by
\[
 R_6=\begin{pmatrix}1&-1\\1&0\end{pmatrix},\qquad
 R_6n_i=n_{i+1},\quad R_6^3=-I,
\]
with subscripts modulo six. Its action on characters is the transpose.
The equality $Q(X-Y,X)=Q(X,Y)$ for $Q=X^2-XY+Y^2$ proves that these are
the same six norm-one vectors as in the integral hexagonal quotient.
Because $R_6$ carries every cone to the next cone, its monomial map and
inverse carry the dual-cone monomial coordinate rings into the corresponding
chart rings. This proves extension across every boundary chart directly,
in agreement with the graph model above.

The source cusp marks its rays, in \S9.3 of the
\href{https://alpo.ge/s6.pdf}{circulated period manuscript}, by
\[
(1,0),(1,-1),(0,-1),(-1,0),(-1,1),(0,1).
\]
The integral matrix $B=\operatorname{diag}(1,-1)$ sends each $n_i$ to the
corresponding source ray, preserves opposite pairs, and is its own inverse.
It induces the exact torus coordinate change $(u,v)\mapsto(u,v^{-1})$.
In that convention the conjugate rotation is
$BR_6B=\left(\begin{smallmatrix}1&1\\-1&0\end{smallmatrix}\right)$ and
the conjugate torus map is $(u,v)\mapsto(uv,u^{-1})$.

This identifies the marked smooth toric normalizations and their six
curves. The source further supplies a normalization map $\nu:dP_6\to W$
identifying each side with its opposite. With that source map retained,
the actual comparison is the composite $S\xrightarrow{f}Y\xrightarrow{B}
dP_6\xrightarrow{\nu}W$. The first map has exactly the three
$\mathbb P^1$ fibres proved above. The last has one preimage off the
boundary, two branches at a general double-curve point, and three at each
triple point, as its local equations $xy=0$ and $xyz=0$ show. The source
marks the two triple-point preimages by the alternating vertex sets
$\{p_{12},p_{34},p_{56}\}$ and $\{p_{23},p_{45},p_{61}\}$.
All side-identification maps remain attached to $\nu$; the fan map alone
does not reconstruct their gluing parameters or identify the singular
quotient with its smooth normalization. No global six-sphere claim is used
in the Picard, torus or contraction proofs.
